% Abstract (~200 words). Numbers: ../FINDINGS.md.
Tools that screen manuscripts for research-integrity problems flag papers that cite retracted work, and retrospective
studies report that such papers are retracted more often. This study tests the signal prospectively on open data.
Predictors were computed on the January 2024 OpenCitations graph using the retractions known at that date, frozen, and
compared with retractions in the Retraction Watch database from February 2024 to October 2026. The design was
pre-specified. Among 10.5 million biomedical papers published 2015--2023, 2.1\% cited a paper already retracted, and
4,323 were retracted later (0.041\%). Citing a retracted paper raised the later retraction rate from 0.035\% to 0.27\%
with one such reference; excluding reviews, it rose to 2.3\% with five or more. Within journal and publication year the odds ratio was 2.99 (95\%
CI 2.73--3.27), and 3.03 when reference-list length was also held fixed, so the signal is not explained by
retraction waves in journals or by long reference lists. As a screen it is weak: 0.9\% of the
1,000 highest-ranked papers were retracted later (2.2\% with reviews set aside), and 84\% of later retractions cite
no retracted paper. When the time
point was moved into the late-2023 retraction wave, the journal's own retraction rate became the far better screen.
NIH-funded papers cited already-retracted work as often as other papers in the same journal and year. Links to
retracted papers are a useful flag but not a screen on their own.
